\chapter{Taxonomy Construction Formalism}\label{app:dag-formalism}\label{sec:meth-construction-summary}

This appendix gives the full architectural and mathematical detail behind the taxonomy
construction pipeline summarised in Section~\ref{sec:arch-taxonomy}
(Chapter~\ref{ch:methodology}).
The pipeline is the instrument rather than the contribution (Section~\ref{sec:approach});
this appendix exists so the exact mechanics are preserved and traceable.

\section{Construction Pipeline Walkthrough}\label{sec:app-dag-pipeline}

\subsection{Seed Initialisation from Ground-Truth Domains}

Before iterative refinement begins, each corpus query is assigned to the anchor node for
its MMLU-Pro ground-truth domain label, creating 14 top-level partitions, each a
\texttt{GraphNode} with initial member set $Q_D^{\text{canonical}}$.
(Anchors sit at depth~1 beneath a formal root at depth~0; depth-conditioned rules below
are stated in implementation depths.)
Anchor immutability is enforced structurally: every destructive topology operation
exempts nodes at depth $\leq 1$, so anchors are never removed, merged, or dissolved.
The anchor's \texttt{sliceDim} is fixed at $d = 256$, the single embedding dimension used
at every depth (Section~\ref{subsec:dim-schedule}).
The ground-truth assignment is initialisation only: trickle routing may reassign a
query, so the adapted partition $Q_D^{\text{adapted}}$ at the end of construction need
not equal $Q_D^{\text{canonical}}$ (Section~\ref{subsec:migration}).

\subsection{Iterative Six-Phase Refinement Loop}\label{tab:pipeline}

After anchor initialisation, \texttt{TaxonomyEngine} executes the six-phase loop already
summarised in Section~\ref{sec:arch-taxonomy} --- embedding
(Section~\ref{sec:meth-mrl}), vMF/NiW fitting (Section~\ref{sec:meth-vmf}), trickle
routing (Section~\ref{sec:meth-routing}), adaptive splitting
(Sections~\ref{sec:meth-split} and~\ref{sec:meth-proposal-gate}), topological
optimisation, and the convergence check (Section~\ref{subsec:quiescence}) --- and is not
re-walked here. Three facts matter downstream.
The loop is global: splitting proceeds per depth across all branches, and queries may
migrate between anchors at every iteration, though routing is skipped entirely in
iteration~1 (Section~\ref{subsec:migration}).
Phase~5 applies its anchor-exempt operations in a fixed, outcome-relevant order ---
starved-leaf handling, sibling fusion, near-duplicate fusion, repeated sole-child
dissolution, and a second starved-leaf pass --- each as a measured proposal under the
gate of Section~\ref{sec:meth-proposal-gate}; each accepted edit re-routes the corpus
and becomes the base state the next proposal is measured against.
Every operator preserves the single-parent invariant
(Section~\ref{app:cross-link-operator}), so each iteration maps a tree to a tree.

\subsection{Cross-Domain Migration}\label{subsec:migration}\label{sec:meth-migration}

At bootstrap, each anchor $D_k$ holds the canonical partition (every query on its
ground-truth anchor with weight $1.0$) with a vMF/NiW self-model fit at $d = 256$.
This assignment is preserved through the first construction iteration: Phase~3
reassignment is skipped in iteration~1, so the first split/optimise/refit pass operates
on the pure ground-truth partition, and geometric reassignment starts at iteration~2,
once real sub-domain structure exists for routing to act on (re-routing immediately in
iteration~1 let queries drift onto merely-similar wrong anchors, and the skip retires the
disabled bias nudge \texttt{enableGtWarmStart}).
From iteration~2 onward, queries are reassigned by the standard trickle routing of
Section~\ref{sec:meth-routing}; \textit{migration} is emergent from this re-routing, and
the \textit{adapted partition} $Q_{D_k}^{\text{adapted}}$ is the anchor-level assignment
after construction completes. Anchor self-models and all sub-domain structure below an
anchor are built from it, and the validation comparison
(Section~\ref{sec:meth-validation}) uses the adapted partition.
The migration rate
$m_{D_k} = |Q_{D_k}^{\text{adapted}} \setminus Q_{D_k}^{\text{canonical}}| /
|Q_{D_k}^{\text{canonical}}|$ is a descriptive diagnostic only: an earlier version made
$\Delta\tau = \tau_{\text{adapted}} - \tau_{\text{canonical}}$ the primary metric of the
second research question, replaced after its premise was tested under pre-registration
and refused (Section~\ref{sec:res-granularity}).

\subsection{Completion, Quiescence, and Reproducibility}\label{subsec:quiescence}

Construction termination is criterion-based with a budget ceiling: the loop runs for at
most \texttt{numIterations} iterations (value in Chapter~\ref{ch:experimental-design}),
and early stopping (\texttt{enableEarlyStopping}, on canonically) is a
\textbf{disjunction of two independent criteria}, evaluated every iteration and reported
jointly in the run log.
The \emph{structural} criterion (\texttt{TaxonomyStabilizer}) counts an iteration
quiescent when the graph edit distance between $\mathcal{G}^{(t)}$ and
$\mathcal{G}^{(t-1)}$ is zero and fires after five consecutive quiescent iterations,
with a floor of $\max(0.8 \cdot |\mathcal{A}|, 5)$ iterations --- at
$|\mathcal{A}| = 14$ anchors, $11$.
The \emph{objective} criterion records
$\Delta J_{\text{edits}} = J_{\text{after edits}} - J_{\text{before edits}}$ and the
re-routing drift $\Delta J_{\text{route}}$, calls an iteration stationary when both lie
within $\tau$, and fires after five consecutive stationary iterations.

\textbf{The canonical run stops at iteration~10, so the structural criterion cannot be
what fired}, its floor being $11$. The run log records \texttt{GED converged=false} at
every iteration and closes with early stopping ``due to convergence (J stationarity)''.
Iterations 8, 9 and 10 are identical on every column of
\texttt{iteration\_metrics.csv} (Section~\ref{sec:res-construction}), which is a property
of the frozen artifact rather than the reason the loop stopped. Stating which arm fired
weakens the guarantee and is recorded for that reason: a structural oscillation leaving
$J$ stationary would satisfy the criterion that stopped this run.
The fixed-point certificate is a separate check on the frozen artifact and has
\textbf{period 1}: it asks only what one further application of the operator would do.
The frozen run's certificate values, quoted in summary form in
Section~\ref{sec:res-construction}, are: edit-phase objective change $0$; re-routing
objective change $-3.33 \times 10^{-16}$; maximum $1 - \cos(\hat{\boldsymbol{\mu}})$
across all nodes $1.11 \times 10^{-16}$; maximum relative change in concentration
$4.25 \times 10^{-13}$; tolerance $10^{-6}$.

\paragraph{Neither criterion observes $\theta$, and the sweep shows what that costs.}
Both criteria read structure and the objective; neither reads the node parameters that
routing depends on. Across the twenty-seed sweep this is not a theoretical gap. Four
seeds stop without certifying, and all four fail on concentration by three to four orders
of magnitude while the structural gate is quiescent. The mechanism in each case is a
\textbf{period-2 cycle}: $\Delta J_{\text{route}}$ alternates in sign at constant
magnitude for every remaining iteration while $\Delta J_{\text{edits}}$ stays exactly
$0$, so the construction is oscillating between two states rather than approaching one.
The certificate, being period~1, correctly refuses; the stopping rule cannot see it.
The distribution across the sweep is bimodal rather than continuous --- sixteen certified
runs sit at a maximum relative $\Delta\kappa$ of $3$--$7 \times 10^{-13}$ and the four
failures at $1.3$--$6.0 \times 10^{-3}$, ten orders apart with nothing between --- so no
run is marginal.
The split by arm is sharper still. Of the three sweep runs that stopped on the
\emph{structural} criterion, three of three failed to certify; of the seventeen that
stopped on objective stationarity, one did (one-sided Fisher exact, $p = 0.0035$).
This is a measured limitation of the stopping rule rather than seed noise, and the fix is
available: requiring the certificate's two $\theta$ terms as a third conjunct, or
refusing the structural arm while $|\Delta J_{\text{route}}| > \tau$, would close it.
Neither is implemented, and the frozen artifact is certified independently of the
question.
A leaf-level log-semantic-volume trajectory is additionally tracked as a convergence
diagnostic and participates in neither criterion.
If neither streak completes, the budget is exhausted and the snapshot is frozen as-is,
with the recorded GED and $J$ trajectories distinguishing quiescent from still-evolving
taxonomies.

Construction is deterministic under a fixed seed: PCA power-iteration initialisation is
seeded from the problem shape, node identifiers are monotonic counters, and
order-sensitive floating-point aggregations iterate their inputs in sorted order, so two
runs at the same seed produce structurally identical taxonomies and cross-configuration
comparisons are exact. Determinism holds \emph{at} a seed and says nothing across seeds:
at the canonical configuration, twenty seeds yield between 72 and 93 leaves against the
frozen artifact's 87 (R3, Section~\ref{sec:arch-requirements}).

\section{Directional Node Model}\label{sec:meth-vmf}

\subsection{von Mises--Fisher Density and MLE}

The vMF distribution on $\mathcal{S}^{d-1}$ with mean direction $\boldsymbol{\mu}$ and
concentration $\kappa \geq 0$ has density
\[
p(\mathbf{x} \mid \boldsymbol{\mu}, \kappa)
= C_d(\kappa) \exp\!\left(\kappa \boldsymbol{\mu}^\top \mathbf{x}\right),
\qquad
C_d(\kappa) = \frac{\kappa^{d/2 - 1}}{(2\pi)^{d/2}\, I_{d/2 - 1}(\kappa)}.
\]
Given $N$ unit-vector observations, the MLE of $\boldsymbol{\mu}$ is the normalised
sample mean $\hat{\boldsymbol{\mu}} = \bar{\mathbf{x}} / \|\bar{\mathbf{x}}\|_2$, and the
Banerjee closed-form approximation~\parencite{banerjee2005vmf}
$\hat{\kappa}_{0} = (\bar{R}\, d - \bar{R}^3)/(1 - \bar{R}^2)$,
$\bar{R} = \|\bar{\mathbf{x}}\|_2$, avoids iterative root-finding for $\kappa$ but
carries an $O(d/N)$ upward bias. It is refined in three steps.
First, it seeds Newton--Raphson steps on the exact likelihood equation
$A_d(\kappa) = \bar{R}$, $A_d(\kappa) = I_{d/2}(\kappa)/I_{d/2-1}(\kappa)$, yielding
$\hat{\kappa}_{\text{ML}}$; at $d = 256$ two or three steps suffice everywhere and the
cap of five is not reached.
Second, the refined estimate is multiplied by a small-sample shrinkage factor
\[
\hat{\kappa}
= \hat{\kappa}_{\text{ML}} \cdot \frac{N - 1}{N + d - 2},
\]
in the spirit of the high-accuracy estimators of Hornik and
Gr\"{u}n~\parencite{hornik2014vmf}.
Because this correction scales with $N$, parents are systematically sharper than their
children --- the reason routing compares \emph{directions} rather than densities
(Section~\ref{sec:meth-routing}).
Third, the shrunk estimate is blended with a parent-anchored empirical-Bayes prior
$\kappa_{\text{parent}}$ (the node's parent's concentration, or a default prior at the
root):
\[
\kappa_{\text{final}}
= (1 - \alpha(\rho))\,\hat{\kappa}
+ \alpha(\rho)\,\kappa_{\text{parent}},
\qquad
\alpha(\rho) =
\begin{cases}
0 & \rho \leq 2,\\
(\rho - 2)/8 & 2 < \rho \leq 10,\\
1 & \rho > 10,
\end{cases}
\]
where $\rho = d/n$ is the dimension-to-effective-sample-size ratio: a node's own data
dominate when it is well supported, and the parent prior takes over as support thins.
At $d = 256$ and the frozen run's $n_{\min} = 55$, a minimum-size leaf sits at
$\rho = 256/55 = 4.65$, hence a blend weight of $\alpha = 0.33$; the fully
prior-dominated $\rho > 10$ regime, reached only below $n = 26$, is not entered by any
leaf of the frozen artifact.
The estimator therefore does not bind leaf size: the floor is chosen on \emph{budget}
grounds and clears the estimator's lower bound near $n = 26$ by a factor of two
(Appendix~\ref{app:tuning-protocol}; Section~\ref{sec:arch-requirements}).
Numerically stable evaluation of $\log I_\nu(\kappa)$ via the Debye asymptotic expansion
is given in Appendix~\ref{app:numerics}.

\subsection{NiW Auxiliary Regulariser}

A Normal-inverse-Wishart (NiW) posterior is fit to each node's member embeddings as an
auxiliary, diagnostic-only summary of its local geometry: routing, migration, and
topology decisions all operate on the vMF geometry (Section~\ref{sec:meth-routing}), and
the stored posterior serves only the leaf log-semantic-volume convergence diagnostic of
Section~\ref{subsec:quiescence}.
The NiW is \textit{not} the conjugate prior for the vMF, whose conjugate is the vMF
itself~\parencite{gopal2014mises}; it is an engineering regulariser whose posterior
predictive is a multivariate Student-$t$, accurate while the tangent-plane approximation
holds. Its hyperparameters, update equations, and predictive density are given in
Appendix~\ref{app:numerics}; because $d = 256$ is fixed at every depth
(Section~\ref{sec:meth-mrl}), a diagonal approximation to $\boldsymbol{\Lambda}_N$ is
stored.
Node centroids are refit directly to the fresh weighted MLE each iteration, subject to a
no-update band that freezes a centroid once its refit direction agrees with the previous
one to within cosine $0.9975$ (an earlier exponential-moving-average smoothing amplified
oscillation in matched A/B runs and was removed).

\section{Fixed-Dimension Matryoshka Embedding Slice}\label{sec:meth-mrl}\label{subsec:dim-schedule}

MRL~\parencite{kusupati2022mrl} trains a single embedding model such that every prefix
slice $\mathbf{x}_{1:m}$ is itself a meaningful embedding
(Section~\ref{sec:taxonomy-construction}).
TaxoArena uses a single fixed slice, $d(\ell) = 256$ for all depths $\ell$, rather than a
per-depth schedule: a larger $d$ increases angular resolution but degrades the
concentration estimate as $d/N$ grows (Section~\ref{sec:meth-vmf}), and a single
dimension keeps all nodes on the same unit hypersphere, so fusion and dissolution
comparisons require no cross-depth down-projection.
Re-normalisation after slicing is mandatory, since MRL guarantees semantic meaningfulness
of $\mathbf{x}_{1:d}$ but not unit norm:
$\mathbf{x}^{(\ell)} = \mathbf{x}_{1:d} / \|\mathbf{x}_{1:d}\|_2$.
The width is a fixed operating point and has not been ablated: it is a compile-time
constant with no configuration path, and an ablation over $d$ would in any case confound
representation width with the concentration estimator's regime and with the
proposal-time PCA geometry (Section~\ref{sec:meth-split}), both of which move with $d$.

\section{Trickle Routing}\label{sec:meth-routing}

Routing separates three decisions that a single absolute floor previously conflated:
whether to descend at all, which siblings compete, and which reached destinations count
as memberships.
At an internal node $v$, the competition set $K(v) = C(v)$ is the node's children.

\textbf{(i) Descent gate.}
The walk descends from $v$ for query $\mathbf{x}$ iff
\[
\max_{c \in K(v)}
\hat{\boldsymbol{\mu}}_c^\top \mathbf{x}
\;\geq\;
(\bar{r}_v - \delta)\,
\hat{\boldsymbol{\mu}}_v^\top \mathbf{x},
\qquad
\bar{r}_v
= \Big\| \sum_{c \in C(v)} \omega_c\, \hat{\boldsymbol{\mu}}_c \Big\|_2
\in [0,1],
\]
where $\omega_c$ are the children's normalised query masses and
$\delta = \texttt{descentMargin} \geq 0$ (canonical $0.12$).
The factor $\bar{r}_v$ makes the bound tight: the parent's direction is approximately
the \emph{normalised} mean of its children's, so the raw parent dot product overstates
what the children can collectively achieve by exactly $1/\bar{r}_v$; the bound is
derived rather than tuned, and at $\delta = 0$ no free parameter enters the decision.
The comparison is between directions, never densities: the $N$-dependent shrinkage of
Section~\ref{sec:meth-vmf} makes parents systematically sharper than children, biasing
cross-level log-density comparisons toward the parent.

A query that fails the gate is neither dropped nor force-assigned: it is held at $v$ with
its full unit weight, so mass is conserved and $\sum_v \sum_q w_v(q) = N$ holds at every
proposal, asserted at runtime.
\textbf{At the canonical $\delta = 0.12$ this outcome never occurs}:
\texttt{residualQueries} is $0$ of $8299$ and the held-out hold rate is $0.00\%$, because
$\delta$ over-provisions the quantity it exists to pay for.
What the Jensen step gives away is the gap between a node's own fitted direction and its
children's mass-weighted resultant; that gap is measured every iteration, and its worst
value is $1 - \cos = 0.0486$ over the frozen artifact's 67 internal nodes and $0.0647$
over all twenty-one runs of this configuration, so the margin of $0.12$ is $1.85\times$
the largest gap ever observed.
Consequently $J$'s cell decomposition (Section~\ref{sec:meth-proposal-gate}) has one cell
per leaf and no others.
This is stated as a limitation rather than a feature: $\delta$ is set well above its
derivation, and while off-configuration runs confirm the mechanism is live (held queries
appear at $\delta = 0.00$ and $0.05$ and none from $0.06$ upward), no run at the frozen
configuration locates the boundary.

\textbf{(ii) Direction-only sibling competition with a cosine beam.}
Admitted siblings are scored
$f_c = \bar{\kappa}\,\hat{\boldsymbol{\mu}}_c^\top \mathbf{x}$ with the \emph{shared
mean} concentration $\bar{\kappa}$ of the level; per-child concentrations and normalisers
are excluded, so a child cannot attract queries through concentration bookkeeping, only
through direction.
The beam keeps every sibling within an additive cosine margin of the best,
\[
\Big\{ c :
\hat{\boldsymbol{\mu}}_c^\top \mathbf{x}
\geq
\max_{c'} \hat{\boldsymbol{\mu}}_{c'}^\top \mathbf{x} - \gamma
\Big\},
\qquad \gamma = \texttt{routingBeamGamma},
\]
the argmax always survives, and per-level transition probabilities renormalise over the
beam.
The walk accumulates log path mass; paths whose absolute posterior falls below
$10^{-30}$ are abandoned as a purely numerical fan-out guard carrying no membership
semantics (an earlier $10^{-4}$ sat inside the range real memberships occupy).

\textbf{(iii) Final membership share.}
After the walk, the query's accumulated masses are normalised over the destinations it
actually reached, and a destination counts as a genuine membership iff it holds at least
\texttt{membershipFloor} of \emph{that query's own} normalised membership.
Because the floor is self-normalised, its meaning is invariant to depth and fan-out; an
absolute product-versus-floor test is not, and mathematically forbids balanced structure
beyond shallow depths.
If no destination clears the share, the single best destination is kept. The floor trims
negligible tails; whether the walk descends at all is the descent gate's decision, not
the floor's.

The \textit{primary} leaf is the argmax of the per-query normalised membership weights
--- a bookkeeping convention for aggregation, not an exclusivity property of the fits.
Held-out queries at evaluation time are routed through the same gate, beam, and share
floor, read-only against the frozen snapshot, capped at \texttt{maxLeafAssignments}
leaves per query as a judge-call-cost bound.
The arena's routing path additionally applies a near-uniform posterior guard (dropping
all secondary memberships when a destination distribution is close to flat) and a second
degree cap, neither of which participates in construction; they are a live candidate
explanation for the routing disagreement of Section~\ref{sec:res-adapted-partition}.
MMLU-Pro labels are never used in the routing decision; they are inspected only
afterward, for the canonical-versus-adapted divergence metrics
(Section~\ref{sec:meth-validation}).

\section{Adaptive Splitting Oracle}\label{sec:meth-split}

Split proposals are generated by vMF-EM clustering; for efficiency the clustering runs in
a PCA-projected subspace whose dimension depends on cluster size (32 dimensions below 100
members, 64 below 500, 128 otherwise), with deterministic, seeded power-iteration
initialisation; the resulting partition is then re-fit and evaluated in the fixed
256-dimensional routing geometry, so no acceptance decision is made in the reduced space.
Only leaves are split. A node is eligible when its mass \emph{and} its effective sample
size $\mathrm{ess} = (\sum_q w_q)^2 / \sum_q w_q^2$ both reach $2 n_{\min}$, and its
depth is below $D_{\max}$; under soft membership these two quantities differ, and the
effective-support test can refuse a node the mass test admits.

A proposed partition $\{S_c\}$ is judged by a \emph{chance-corrected separation score}.
For a set $S$ of unit vectors, define the total pairwise cosine dissimilarity
\[
W(S) = |S|^2 - \Big\|\sum_{\mathbf{x} \in S} \mathbf{x}\Big\|_2^2
= \sum_{i \neq j \in S} \left(1 - \mathbf{x}_i^\top \mathbf{x}_j\right),
\]
computable in $O(|S|\,d)$ from the member sum alone.
The observed within-cluster scatter of the partition is $\sum_c W(S_c)$; under a
uniformly random partition of the same $n$ points into the same cluster sizes $\{n_c\}$,
its exact expectation is
\[
\mathbb{E}\Big[\sum_c W(S_c)\Big]
= W(S_{\text{total}}) \cdot
\sum_c \frac{n_c (n_c - 1)}{n (n - 1)},
\]
because each unordered pair lands inside cluster $c$ with probability
$n_c(n_c-1)/(n(n-1))$.
The score is
\[
s = 1 - \frac{\sum_c W(S_c)}
{\mathbb{E}\left[\sum_c W(S_c)\right]},
\]
which equals $1$ for a perfect partition (each cluster internally identical), is
approximately $0$ for a random partition, negative for an anti-clustered one, and exactly
$0$ for the degenerate $k = 1$ ``partition'', so a non-separating wrapper can never clear
a positive threshold.
The chance correction removes the mechanical dependence a raw within/total ratio has on
$k$ and on the cluster-size profile, so a single threshold means the same thing
everywhere in the graph.
This score is \emph{not} a Dasgupta cost delta; Dasgupta's cost
\parencite{dasgupta2016cost} is retained only as the whole-tree Total Dasgupta Cost
evaluation metric (Appendix~\ref{app:metrics}).

Acceptance proceeds in four stages, of which the EM clustering is only the first.
\begin{enumerate}
\item \textbf{Proposal.} vMF-EM at a \emph{forced} $k$, tried in ascending order
$k = 2, 3, \ldots, \texttt{maxK}$, with the first $k$ whose proposal survives every later
stage winning and no further $k$ attempted. Selection is by first acceptance rather than
by comparing candidates, which keeps the split bit-identical at a fixed seed: no float
comparison decides between two $k$ values. \texttt{maxK} $= 4$ is a cost bound.
\item \textbf{Routed sustainability.} Proposal vMF components are refit at $d = 256$ and
every target query is re-assigned by the same level-local posterior the trickler uses;
the split is rejected unless \emph{every} child holds at least $n_{\min}$ members under
that routed assignment.
This closes the loop between birth and survival --- a split occurs only if routing will
sustain every child it creates --- and eliminated a persistent spawn-and-starve churn in
which the majority of pruned nodes died within two iterations of birth. Children are then
refit on their routed populations.
\item \textbf{Min-pairwise separation.} Every \emph{pair} of routed children must score
$s \geq \epsilon_{\text{sep}}$ on the chance-corrected statistic above.
The gate is pairwise rather than $k$-way for gate consistency: the same statistic, at
the same bar, is what the sibling merger applies to \emph{destroy} a pair, and a joint
$k$-way score of $0.06$ can contain a pair at $0.015$ that the merger fuses in the same
iteration --- the split/fuse limit cycle. With one statistic at one bar the two
acceptance regions are disjoint, so nothing creatable is immediately destroyable.
It is also the binding gate empirically: $129$ of the frozen run's $373$ proposal records
are min-pairwise rejections, against $5682$ min-pairwise rejections and $0$ joint $k$-way
rejections across every run in the repository; the joint score survives only as the
per-node \texttt{dasguptaDeltaNorm} diagnostic that the within-node null reads.
\item \textbf{Global gate.} The split is committed only if it clears the standard-error
rule of Section~\ref{sec:meth-proposal-gate}. Local quality is already established by
stages 1--3; the global test asks the different question of whether the improvement is
larger than its own measurement error.
\end{enumerate}
The recursion depth is hard-capped at $D_{\max}$, which the frozen artifact does not
reach ($D_{\max} = 8$ against a realised depth of 6).
The threshold $\epsilon_{\text{sep}}$ (\texttt{proposalSeparationBar}) has no closed-form
optimum. Its value in the frozen run is $\epsilon_{\text{sep}} = 0.025$, positioned
between two measured nulls: an isotropic null whose $p_{95}$ spans $0.0055$--$0.0093$
over $n = 75\ldots900$, and the within-node null of a node's own elongation, whose median
lies above it. The bar is therefore conservative by a factor of $2.7$--$4.5$ against
``separated at all'' and permissive against ``separated by more than this node's own
elongation''; the second is the source of limitation~1 in Section~\ref{sec:arch-limits}.
The \emph{same} score at the \emph{same} bar also gates sibling fusion, which is the
property stage~3 depends on. A single constant serving both a \emph{level} and a
\emph{difference} would be the threshold-reuse defect of rule~D4
(Section~\ref{sec:meth-discipline}); the two roles are decoupled into independent
configuration fields, and the difference role is not used, $k$ being chosen by first
acceptance rather than marginal improvement.

\section{The Proposal Gate}\label{sec:meth-proposal-gate}

Every structural edit (split, fusion, prune, and dissolution) is applied as a
\emph{measured proposal} rather than by its local rule alone.
The gate operates on a global objective $J$: the chance-corrected separation score of
Section~\ref{sec:meth-split} lifted to the whole corpus, with one cell per leaf, weighted
by soft membership.
The cell decomposition is exhaustive because routing is: at the canonical descent
margin every query reaches a leaf (Section~\ref{sec:meth-routing}), so no cell is empty.
The proposal procedure is: snapshot the current state (topology, node parameters, and
assignments); apply the candidate edit tentatively; recompute descent-gate statistics;
clear and fully re-route every query through the modified taxonomy; measure
$J' = J(\mathcal{G}')$; and commit under an acceptance rule expressed in units of the
\emph{paired} standard error of $\Delta J$ itself:
\[
\text{accept}
\iff
\begin{cases}
\Delta J > \max\bigl(\tau,\; z^{*}\cdot
\widehat{\mathrm{SE}}_{\text{paired}}(\Delta J)\bigr)
& \widehat{\mathrm{SE}}_{\text{paired}}(\Delta J) > 0,\\[4pt]
\Delta |V| < 0
& \widehat{\mathrm{SE}}_{\text{paired}}(\Delta J) = 0,
\end{cases}
\qquad
\Delta J = J' - J,
\qquad z^{*} = 2.0 .
\]
The test is one-sided; on rejection the snapshot is restored.
$\widehat{\mathrm{SE}}_{\text{paired}}(\Delta J)$ is a paired Dirichlet bootstrap over
200 replicates at a fixed seed, resampling only queries whose cell assignment changed
between the two states, so both sides share a corpus draw and the shared variance
cancels.

The second branch is not a fallback tolerance but a different case:
$\widehat{\mathrm{SE}}_{\text{paired}} = 0$ identifies a pure structural edit that moves
no query --- a pass-through dissolution --- for which $\Delta J$ is exactly $0$ and $z$
is undefined; the size test lets those commit. The two branches are alternatives rather
than layers: where the standard error is positive there is no lexicographic tie-break, so
a $J$-neutral simplification cannot commit. Exactly one of the frozen run's 70 accepted
edits took the second branch.
The floor $\tau$ (\texttt{tau}, canonical $10^{-6}$) is what makes a termination bound
available: every accepted edit raises $J$ by at least $\tau$ and $J$ is bounded above by
$1$, so at most $(1 - J_0)/\tau$ edits can ever commit; with $z^{*}\cdot\mathrm{SE}$
alone a vanishing standard error would admit a vanishing improvement and that bound
collapses.

\textbf{Why the rule is expressed this way is the argument, and it is separate from any
claim that it improves the tree.}
The size of a defensible improvement is not a property of the corpus but of the edit:
across the frozen run's 69 accepted edits that move queries, $\mathrm{SE}(\Delta J)$
spans $13.4\times$, from $3.27\times10^{-5}$ to $4.38\times10^{-4}$, so no single
absolute tolerance is correct at both ends of that range.
The gate declined eight split sites whose $\Delta J$ exceeded $\tau$, four at $z < 1$
and the worst at $z = 0.089$ ($\Delta J = 9.46\times10^{-6}$ against
$\mathrm{SE}(\Delta J) = 1.06\times10^{-4}$); a $\tau$-only rule commits all eight, and
committing an edit whose measured improvement is an eleventh of its own measurement error
is incoherent regardless of what it does downstream.
In the frozen run the $\tau$ arm never bound: the smallest operative threshold across
the 173 gate rejections is $8.83\times10^{-5}$ ($88\tau$) and the smallest
$z^{*}\cdot\mathrm{SE}$ among accepted edits $6.54\times10^{-5}$ ($65\tau$), so the rule
reduced in practice to $\Delta J > 2\,\mathrm{SE}$ at every site and the threshold that
governed every decision was measured rather than set.
The gate's two natural empirical corroborations --- that it tightens run-to-run
stability \emph{across seeds}, and that it raises $J$ --- were both tested and both
refuted (Section~\ref{sec:res-construction}), and it is adopted on coherence alone.

\paragraph{The neutrality measurements.} Demoted from
Section~\ref{sec:res-construction}, which keeps the verdicts. The
gate's effect on run-to-run structure is an additive offset of about
$4.4$ leaves. The dispersion reading needs its mechanics stated,
because the raw numbers point the wrong way: the coefficient of
variation of leaf counts rose from $14.3\%$ to $15.7\%$ under the gate
and was initially read as a degradation, but the standard deviation is
unchanged (variance ratio $1.05$ against an $F(4,4)$ critical value of
$6.39$; the correlation between a seed's leaf-count reduction and its
deviation from the mean is $+0.006$), and the coefficient of variation
rises mechanically because the mean fell. The correct statement is
that no dispersion change is detectable in either direction, the
second instance of rule~D3 (Section~\ref{sec:meth-discipline}). On the
objective: $J$ is higher on 2 of 5 seeds, exact two-sided sign test
$p = 1.000$, median $\Delta J = -0.0001$; the apparent gain was seed
42 alone. An earlier draft quoted $p = 0.812$ for this comparison,
which no tail of a sign test on five seeds can return; it is
withdrawn. Caveat: the five-seed sweep's leaf counts predate a
correction to the splitter's routing check and are \textsc{void as to
values} (marker defined in Section~\ref{sec:res-construction}); the
two negative results are comparisons internal to the sweep and are
unaffected, but the underlying counts need re-deriving.

The same argument rules out selecting the branching factor $k$ by $\arg\max \Delta J$
over candidates: a maximum over noisy estimates is biased upward even when all
candidates are equivalent, so the $k$-fallback takes the \emph{lowest} $k$ among
candidates that each independently cleared the same gate.
Setting $z^{*} = 0$ recovers the legacy lexicographic rule
$\Delta J > \tau \;\vee\; (|\Delta J| \leq \tau \wedge \Delta|V| < 0)$, retained as an
attributable baseline arm. The rule is uniform across edit types, making accept and
reject decisions comparable between splits and shrink edits; every proposal writes
$\Delta J$, its standard error, the resulting $z$, and the binding threshold to the
proposal log.

The gate has two structural consequences.
Sole-child dissolution needs no separate justification: a parent's only tree child
separates nothing (the score of Section~\ref{sec:meth-split} is identically $0$ at
$k = 1$), and the gate confirms empirically that removing it does not degrade $J$;
top-level nodes are exempt.
Starved-leaf handling (a leaf whose branch falls below $n_{\min}$, the same floor that
gates birth) is a three-way choice --- keep, absorb into the parent, or fuse into the
nearest sibling --- decided by whichever option measures best under the gate.

\paragraph{What the coarsening half actually contributed.}
The gate applies uniformly to growth and coarsening, but the frozen run used the two
halves asymmetrically: of 70 accepted edits, 68 were splits and \textbf{2} were
coarsening moves --- both sole-child dissolutions, both in iteration~2. No fusion was
ever accepted, none was proposed after iteration~3, and the remaining ten coarsening
proposals were declined at $z$ between $-7.2$ and $-14.6$, strongly $J$-degrading rather
than marginal. The design's growth/coarsening symmetry is therefore stated but not
exercised by this artifact: what shaped the frozen tree was the split gate declining
candidates, not the merger reclaiming them.

\section{The Single-Parent Invariant}\label{app:cross-link-operator}

Every topology operator maps a tree to a tree, and one of them needs an explicit
restriction to do so: fusing two nodes under different parents would leave the survivor
holding both parent edges, producing a multi-parent node, so fusion proposals are
restricted to pairs sharing a parent.
The structural-integrity gates of Appendix~\ref{app:tuning-protocol} check the invariant
directly by counting multi-parent nodes, which must be zero; in the frozen artifact all
154 nodes have exactly one parent and \texttt{crossDomainNodes} is $0$.
An operator that deliberately created second parents was built, measured, and retired
before this construction was frozen. It is reported as a measured negative result in
Section~\ref{sec:poly-negative} and forms no part of the method described here.

\paragraph{Evaluation-time multi-leaf assignment} of held-out queries is a separate
mechanism. It applies the same routing read-only to the frozen snapshot, capped at
\texttt{maxLeafAssignments}, and does not modify the topology.
It does \emph{not} determine which leaves a query contributes comparisons to: for every
arena run reported in this thesis the propagation path received each query's primary leaf
only, and the secondary memberships were computed and discarded
(Appendix~\ref{app:numerics}).

\paragraph{The admission record and the provenance caveat.} Demoted
from Section~\ref{sec:poly-negative}, which keeps the retiring
evidence. That admission was decided by evaluation order rather than
by evidence is explicit in the proposal sequence: one node was
accepted at four hosts and rejected at two, and acquired five parents
by iteration~1 (median capture fraction $f$ of $0.802$ accepted
against $0.818$ rejected). One caveat on provenance bounds what those
numbers can be checked against: they come from a construction the
current source can no longer produce --- the operator was deleted, and
with it the proposal type its records were written under, so the
proposal logs of the frozen artifact and of every run that postdates
the deletion contain no cross-link rows. The measurements stand as a
record of why the operator was retired and cannot be re-derived from
any artifact this thesis reports.
